\section{Connected Data}



\subsection{Update Protocols}

The recorded lectures are provided \myhref{https://www.youtube.com/watch?v=-P4Qks9J3o0}{here (Introduction)} and \myhref{https://www.youtube.com/watch?v=B70b-uyNHpw}{here (Update Protocols)}.

\begin{ejercicio}
    What is an update protocol and how does it differ from stream compression?
\end{ejercicio}

\begin{ejercicio}
    Explain the update protocol ``Periodic Reporting'' and describe a situation in which it works very well and a situation in which it does not work well. Which other update protocol can we use to solve the problem?
\end{ejercicio}

\newpage\subsection{Data Models}

The recorded lecture is provided \myhref{https://www.youtube.com/watch?v=ZZ3H8QTZKv8}{here}.

\begin{ejercicio}
    Give an example for an information that a Key-Value Model cannot represent but a Linked-data Based Model can.
\end{ejercicio}

\newpage\subsection{Stream Processing and Complex Event Processing}

The recorded lectures are provided \myhref{https://www.youtube.com/watch?v=lAdIGfW8fUg&feature=youtu.be}{here (Introduction)}, \myhref{https://www.youtube.com/watch?v=3oDF0BbPR0A&feature=youtu.be}{here (Stream Processing)} and \myhref{https://www.youtube.com/watch?v=8hbLKM3zH9w}{here (Complex Event Processing)}.
\begin{ejercicio}
    Compare Stream Processing with Complex Event Processing by discussing the different information models, what they each focus on and where they come from.
\end{ejercicio}

\begin{ejercicio}
    What is a continuous query? Give an example for one in pseudo code (i.e., define your own language for continuous queries and provide an example with it). Include typical features like, e.g., a window(-ing) concept into your language.
\end{ejercicio}

\begin{ejercicio}
    What are Event-Condition-Action rules and how do they differ from continuous queries? What is similar?
\end{ejercicio}

\newpage\subsection{IoT Machine Learning}

The recorded lecture is provided \myhref{https://www.youtube.com/watch?v=S4ycy_yPp2c&feature=youtu.be}{here}.

\begin{ejercicio}
    Why would we want to move machine learning to the edge or even into the smart IoT thing? What are potential advantages? What are challenges and when would you not move it but keep it in the Cloud?
\end{ejercicio}

\begin{ejercicio}
    Briefly explain how magnitude-based pruning is performed.
\end{ejercicio}

\begin{ejercicio}
    Explain the difference between pruning and quantization of an artificial neural network. What is binarization and why has it the potential to be implemented extremely efficiently in hardware?
\end{ejercicio}